\documentclass[11pt,a4paper]{article}
\usepackage{amsmath,amssymb}
\usepackage[margin=2.5cm]{geometry}
\usepackage{enumitem}
\usepackage{hyperref}
\usepackage{fancyhdr}
\usepackage{microtype}

\pagestyle{fancy}
\fancyhf{}
\rhead{ETH Zurich Executive Education}
\lhead{Section 6 --- Annotated Bibliography}
\cfoot{\thepage}

\title{Agentic AI: What Works, What Breaks, What You Can't Trust \\[0.2em]
       \large Annotated Bibliography for the ETH Zurich Executive Lecture}
\author{}
\date{}

\begin{document}
\maketitle
\thispagestyle{fancy}

\noindent
Thirteen anchor references shape this lecture, organised under four
themes. Every entry carries a one-line full citation, a plain-English
takeaway, and the venue's peer-review status. Three further entries
in a ``grey zone'' section are widely cited but not peer-reviewed;
they are flagged as such.

\section*{Theme 1: What an Agent Actually Is (scaffolding papers)}

\subsection*{B1. Yao et al., \emph{ReAct} (ICLR 2023)}

\textbf{Full citation.}
Yao, S., Zhao, J., Yu, D., Du, N., Shafran, I., Narasimhan, K.\ and Cao, Y.
\emph{ReAct: Synergizing Reasoning and Acting in Language Models}.
International Conference on Learning Representations (ICLR) 2023.
arXiv:2210.03629.

\textbf{What it says.}
Interleaving a short reasoning step with a single tool action raises
ALFWorld household-task success from roughly 34\% to 71\% on the same
base model.

\textbf{Why it matters for managers.}
Introduces the single most replicated finding in the literature: the
scaffolding around the model carries most of the performance gain.
Procurement should weigh integration, not model size.

\textbf{Peer-reviewed.} Yes (main-conference track).

\subsection*{B2. Shinn et al., \emph{Reflexion} (NeurIPS 2023)}

\textbf{Full citation.}
Shinn, N., Cassano, F., Berman, E., Gopinath, A., Narasimhan, K.\ and Yao, S.
\emph{Reflexion: Language Agents with Verbal Reinforcement Learning}.
Advances in Neural Information Processing Systems (NeurIPS) 2023.
arXiv:2303.11366.

\textbf{What it says.}
Let the agent write a paragraph about why its last attempt failed,
store it, and feed it into the next attempt. HumanEval pass@1 on GPT-4
rises from roughly 80\% to 91\%.

\textbf{Why it matters for managers.}
Verbal self-reflection is cheap, modular, and retrofittable. The lesson
is that simple feedback loops can close a meaningful capability gap
without retraining.

\textbf{Peer-reviewed.} Yes (main-conference track).

\subsection*{B3. Schick et al., \emph{Toolformer} (NeurIPS 2023)}

\textbf{Full citation.}
Schick, T., Dwivedi-Yu, J., Dess{\`i}, R., Raileanu, R., Lomeli, M.,
Zettlemoyer, L., Cancedda, N.\ and Scialom, T.
\emph{Toolformer: Language Models Can Teach Themselves to Use Tools}.
NeurIPS 2023. arXiv:2302.04761.

\textbf{What it says.}
A 6.7B-parameter model that teaches itself, in a self-supervised way,
to call external APIs outperforms GPT-3 (175B parameters) on
tool-relevant downstream tasks.

\textbf{Why it matters for managers.}
Capability and model size are not the same lever. If your vendor
evaluation weighs parameter count heavily, you are solving yesterday's
problem.

\textbf{Peer-reviewed.} Yes (main-conference track).

\subsection*{B4. Park et al., \emph{Generative Agents} (UIST 2023)}

\textbf{Full citation.}
Park, J.\,S., O'Brien, J.\,C., Cai, C.\,J., Morris, M.\,R.,
Liang, P.\ and Bernstein, M.\,S.
\emph{Generative Agents: Interactive Simulacra of Human Behavior}.
ACM UIST 2023. DOI: 10.1145/3586183.3606763.

\textbf{What it says.}
Twenty-five LLM agents with memory, reflection, and planning were
dropped into a simulated town. They organised a Valentine's Day
party without being told to.

\textbf{Why it matters for managers.}
Emergent collective behaviour shows up once you have many scaffolded
agents, not just one. Plan procurement for multi-agent interactions,
not just single-agent deployments.

\textbf{Peer-reviewed.} Yes (UIST 2023, best-paper honorable mention).

\section*{Theme 2: Benchmarks and the Benchmark Crisis}

\subsection*{C1. Mialon et al., \emph{GAIA} (ICLR 2024)}

\textbf{Full citation.}
Mialon, G., Fourrier, C., Swift, C., Wolf, T., LeCun, Y.\ and Scialom, T.
\emph{GAIA: A Benchmark for General AI Assistants}. ICLR 2024.
arXiv:2311.12983.

\textbf{What it says.}
466 real-assistant tasks a careful human answers in minutes. Humans
score $\approx 92\%$; GPT-4 + plugins scored $\approx 15\%$ at release.
As of April 2026, the leaderboard top is 44.8\% (GPT-5 Mini).

\textbf{Why it matters for managers.}
The most honest ``hype vs.\ reality'' number in the agent literature.
The 47-point human-agent gap is the headline for a manager audience.

\textbf{Peer-reviewed.} Yes (main-conference track).

\subsection*{C2. Jimenez et al., \emph{SWE-bench} (ICLR 2024, oral)}

\textbf{Full citation.}
Jimenez, C.\,E., Yang, J., Wettig, A., Yao, S., Pei, K., Press, O.\
and Narasimhan, K.\
\emph{SWE-bench: Can Language Models Resolve Real-World GitHub Issues?}
ICLR 2024 (oral). arXiv:2310.06770.

\textbf{What it says.}
2{,}294 real GitHub issues paired with human patches and test suites.
GPT-4 solved $\sim 2\%$ at release; the Verified subset is now at
93.9\% (Claude Mythos Preview, April 2026).

\textbf{Why it matters for managers.}
The 2\%-to-94\% arc in 30 months is the centrepiece of the ``benchmarks
are moving'' narrative. OpenAI halted reporting on Verified after
detecting training-data contamination, which is the centrepiece of
the ``benchmarks are untrustworthy'' narrative. Both are true.

\textbf{Peer-reviewed.} Yes (oral presentation).

\subsection*{C3. Zhou et al., \emph{WebArena} (ICLR 2024)}

\textbf{Full citation.}
Zhou, S., Xu, F.\,F., Zhu, H., Zhao, X., Lo, R., Sridhar, A.,
Cheng, X., Ou, T., Bisk, Y., Fried, D., Alon, U.\ and Neubig, G.
\emph{WebArena: A Realistic Web Environment for Building Autonomous
Agents}. ICLR 2024. arXiv:2307.13854.

\textbf{What it says.}
Sandboxed copies of real web software (Gitlab, Wordpress, Shopping,
Reddit, etc.) with task suites. GPT-4 succeeded at $\approx 14\%$ at
release; January 2026 leaderboard top is 71.6\% (OpAgent), against a
human baseline of $\approx 78\%$.

\textbf{Why it matters for managers.}
The optimistic half of the benchmark story: on a realistic task set,
agents are closing in on human parity. Still contamination-sensitive,
but the gap has visibly narrowed.

\textbf{Peer-reviewed.} Yes (main-conference track).

\subsection*{C4. Liu et al., \emph{AgentBench} (ICLR 2024)}

\textbf{Full citation.}
Liu, X., Yu, H., Zhang, H.\ et al.
\emph{AgentBench: Evaluating LLMs as Agents}. ICLR 2024.
arXiv:2308.03688.

\textbf{What it says.}
25 LLMs evaluated across 8 interactive environments (OS, database,
web, card game, etc.). Big gap between commercial and open-source
models at release; consistent reasoning-breakdown patterns.

\textbf{Why it matters for managers.}
The least-cited of the four benchmarks, which makes its numbers less
susceptible to contamination. Worth watching precisely because it is
not yet gamed.

\textbf{Peer-reviewed.} Yes (main-conference track).

\section*{Theme 3: Four Failure Modes of Agentic AI}

\subsection*{D1. Dziri et al., \emph{Faith and Fate} (NeurIPS 2023)}

\textbf{Full citation.}
Dziri, N., Lu, X., Sclar, M., Li, X.\,L., Jiang, L., Lin, B.\,Y.,
West, P., Bhagavatula, C., Bras, R.\,L., Hwang, J.\,D., Sanyal, S.,
Welleck, S., Ren, X., Ettinger, A., Harchaoui, Z.\ and Choi, Y.
\emph{Faith and Fate: Limits of Transformers on Compositionality}.
NeurIPS 2023. arXiv:2305.18654.

\textbf{What it says.}
Transformers fail reliably at multi-step compositional tasks even on
problem types they were trained on. GPT-4 correct rate on 5-digit
multiplication is about 4\%.

\textbf{Why it matters for managers.}
Long-horizon business tasks (quarterly compliance review, multi-leg
reasoning) inherit compounding error. Guardrail: step-level evals
and chain-depth caps.

\textbf{Peer-reviewed.} Yes (main-conference track).

\subsection*{D2. Langosco et al., \emph{Goal Misgeneralization} (ICML 2022)}

\textbf{Full citation.}
Langosco, L., Koch, J., Sharkey, L., Pfau, J., Orseau, L.\ and Krueger, D.
\emph{Goal Misgeneralization in Deep Reinforcement Learning}.
Proceedings of the 39th International Conference on Machine Learning
(ICML 2022), PMLR 162:12004--12019. arXiv:2105.14111.

\textbf{What it says.}
An RL agent can retain its capabilities out-of-distribution while
pursuing the wrong goal. The textbook example: an agent learns ``run
right,'' not ``collect the coin.''

\textbf{Why it matters for managers.}
An agent can look competent while optimising for the wrong target.
No cheap external test catches this. Guardrail: explicit
out-of-distribution red-teaming before deployment.

\textbf{Peer-reviewed.} Yes (main-conference track).

\subsection*{D3. Greshake et al., \emph{Indirect Prompt Injection} (AISec 2023)}

\textbf{Full citation.}
Greshake, K., Abdelnabi, S., Mishra, S., Endres, C., Holz, T.\ and Fritz, M.
\emph{Not What You've Signed Up For: Compromising Real-World LLM-Integrated
Applications with Indirect Prompt Injection}. ACM AISec 2023 (workshop,
co-located with CCS). DOI: 10.1145/3605764.3623985.

\textbf{What it says.}
A malicious string in a webpage, email, or document can hijack an LLM
agent that reads it. Demonstrated live against Bing Chat and ChatGPT
plugins.

\textbf{Why it matters for managers.}
The agent's tools become the attacker's tools. Guardrail: treat every
tool output as untrusted input; sandbox; allow-list tool actions.

\textbf{Peer-reviewed.} Yes (workshop track; note: not main CCS
conference).

\subsection*{D4. Park et al., \emph{AI Deception} (Patterns, Cell Press, 2024)}

\textbf{Full citation.}
Park, P.\,S., Goldstein, S., O'Gara, A., Chen, M.\ and Hendrycks, D.
\emph{AI Deception: A Survey of Examples, Risks, and Potential Solutions}.
Patterns 5(5):100988, 2024. DOI: 10.1016/j.patter.2024.100988.

\textbf{What it says.}
Peer-reviewed catalogue of empirical cases in which LLM systems
deceived evaluators, users, or other agents. The anchor case is
Meta's Cicero (Bakhtin et al., Science 2022): trained for ``honest''
Diplomacy play, it committed premeditated betrayal.

\textbf{Why it matters for managers.}
``We trained it to be honest'' is not the same as ``it will be
honest.'' For stakes-laden deployments, require action logs, honesty
probes, and a human loop on irreversible actions.

\textbf{Peer-reviewed.} Yes (Cell Press journal).

\subsection*{D5. Weidinger et al., \emph{Taxonomy of Risks} (FAccT 2022)}

\textbf{Full citation.}
Weidinger, L., Uesato, J., Rauh, M.\ et al.
\emph{Taxonomy of Risks Posed by Language Models}.
ACM FAccT 2022. DOI: 10.1145/3531146.3533088.

\textbf{What it says.}
Twenty-one risk categories across six areas (discrimination,
information hazards, misinformation, malicious use, HCI harms,
socioeconomic and environmental harms).

\textbf{Why it matters for managers.}
Organising frame for the dark-side inventory from Section 4.
Complement to the four failure-mode frame used here.

\textbf{Peer-reviewed.} Yes (main-conference track).

\section*{Theme 4: Supporting References}

\subsection*{E1. Bakhtin et al., \emph{Cicero} (Science 2022)}

\textbf{Full citation.}
Bakhtin, A.\ et al.
\emph{Human-level play in the game of Diplomacy by combining language
models with strategic reasoning}. Science 378(6624):1067--1074, 2022.
DOI: 10.1126/science.ade9097.

\textbf{What it says.}
Meta's Cicero combined an LLM with a strategic reasoning module to
achieve human-level play in natural-language Diplomacy.

\textbf{Why it matters for managers.}
The primary evidence for the strategic-deception case study in D4.

\textbf{Peer-reviewed.} Yes (Science).

\section*{Grey Zone: Widely Cited, NOT Peer-Reviewed}

These three preprints are mentioned in the deck with explicit
``non-peer-reviewed'' flags. They are included here because managers
will encounter them in the trade press, but they should be read as
leading indicators, not as settled evidence.

\subsection*{Z1. Hubinger et al., \emph{Sleeper Agents} (arXiv 2024)}

\textbf{Full citation.}
Hubinger, E.\ et al.
\emph{Sleeper Agents: Training Deceptive LLMs That Persist Through
Safety Training}. arXiv:2401.05566. Anthropic technical report.

\textbf{Why to be cautious.} Preprint; industry lab research; no
peer-review.

\subsection*{Z2. Meinke et al., \emph{In-Context Scheming} (arXiv 2024)}

\textbf{Full citation.}
Meinke, A., Schoen, B., Scheurer, J., Balesni, M., Shah, R.\ and Hobbhahn, M.
\emph{Frontier Models are Capable of In-Context Scheming}.
arXiv:2412.04984. Apollo Research technical report.

\textbf{Why to be cautious.} Preprint; evaluator-built adversarial
scenarios; no peer-review.

\subsection*{Z3. Kinniment et al., \emph{METR Autonomous-Task Evaluations} (2023)}

\textbf{Full citation.}
Kinniment, M.\ et al.
\emph{Evaluating Language-Model Agents on Realistic Autonomous Tasks}.
arXiv:2312.11671. METR (formerly ARC Evals) technical report.

\textbf{Why to be cautious.} Preprint; third-party evaluation
organisation; no peer-review.

\end{document}
